\documentclass[12pt]{article}
\usepackage[utf8]{inputenc}
\usepackage{graphicx}
\usepackage{animate}
%导入一些用到的宏包
\usepackage{amsmath,bm,amsfonts,amssymb,enumerate,epsfig,bbm,calc,color,ifthen,capt-of,multimedia,bm,multirow}
\usepackage{blindtext}
\usepackage{booktabs}
\usepackage{MnSymbol}
\usepackage{url}
\usepackage{natbib}
\usepackage{changepage}
\usepackage{fancyhdr}
\usepackage[flushleft]{threeparttable}
\newtheorem{theorem}{Theorem}
\usepackage{setspace}
\onehalfspacing

\linespread{1.5}
\textwidth=39pc
\textheight=50pc
\oddsidemargin=0pc
\evensidemargin=0pc
\headsep=0pt
%\headheight=.2cm
\topmargin=0pc
\parindent=1.7pc
\parskip=0pt
\usepackage[flushleft]{threeparttable}


%\newtheorem{theorem}{Theorem}
\DeclareMathOperator*{\argmax}{arg\,max}
\DeclareMathOperator*{\argmin}{arg\,min}


\title{Weighted Q-learning for optimal dynamic treatment regimes with nonignorable missing covariates}
\author{Jian Sun \thanks{School of Data Science,  Fudan University, Shanghai 200433, China. Email: jsun19@fudan.edu.cn.}, Bo Fu \thanks{School of Data Science,  Fudan University, Shanghai 200433, China. Email: fu@fudan.edu.cn.}, and Li Su \thanks{MRC Biostatistics Unit, University of Cambridge, Cambridge, CB2 0SR, UK. Email: li.su@mrc-bsu.cam.ac.uk.}}


\date{}


\providecommand{\keywords}[1]{\textbf{\textit{Keywords  }} #1}

\begin{document}

\maketitle

\begin{abstract}
Dynamic treatment regimes (DTRs)  formalize medical decision-making as a sequence of rules for different stages, mapping patient-level information to recommended treatments. In practice, estimating an optimal DTR using observational data from electronic medical  record (EMR) databases can be complicated by nonignorable missing covariates resulting from %that are \iffalse missing not at random (MNAR) \fi  due to 
informative monitoring of patients. Since complete case analysis can provide consistent estimation of outcome model parameters under the assumption of outcome-independent missingness\iffalse \citep{Yang_Wang_Ding_2019}\fi, Q-learning is a natural approach to accommodating nonignorable missing covariates. 
%Building on the consistency of the outcome model estimated by the complete case analysis method under the assumption of future-independent missingness at the last stage, we consider Q-learning to estimate dynamic treatment regimes. 
However, the backward induction algorithm used in Q-learning can introduce challenges, as nonignorable missing covariates at later stages can result in nonignorable missing pseudo-outcomes at earlier stages, leading to suboptimal DTRs, even if the longitudinal outcome variables are fully observed. To address this unique missing data problem in DTR settings, we propose two weighted Q-learning approaches where inverse probability weights for missingness of the pseudo-outcomes are obtained through estimating equations with valid nonresponse instrumental variables or sensitivity analysis. The asymptotic properties of the weighted Q-learning estimators are derived, and
the finite-sample performance of the proposed methods is evaluated and compared with alternative methods through extensive simulation studies. Using EMR data from the Medical Information Mart for Intensive Care database, we apply the proposed methods to investigate the optimal fluid strategy for sepsis patients in intensive care units.
\end{abstract}

\begin{keywords}
Backward-induction-induced missing pseudo-outcome; Future-independent missingness; Nonignorable missing data; Nonresponse instrumental variable; Q-learning; Sensitivity analysis. 
\end{keywords}




\section{Introduction}

Dynamic treatment regimes (DTRs) formalize medical decision-making as a sequence of decision rules, each corresponding to a specific decision stage, that map available information at patient-level to a recommended treatment \citep{Murphy_2003}. Optimal DTRs have been estimated and evaluated in various biomedical areas, such as cancer \citep{wang2012evaluation} and sepsis \citep{Yu_Zhang_2022}. In practice, estimating optimal DTRs becomes challenging when there are missing data in covariates, treatments, or outcomes.

There is a rich literature on identifying optimal DTRs with complete data. Various methods have been developed, including  Q-learning \citep{Murphy_2003}, A-learning \citep{Schulte_et_al_2014}, and their numerous variants \citep{Kosorok_Laber_2019}. These methods rely on approximate dynamic programming methods. Direct-optimization methods, such as simultaneous outcome weighted learning \citep{Zhao_et_al_2015}, and model-based planning via g-computation \citep{Xu_et_al_2016} have also been employed. In contrast, limited research to date has been devoted to accommodating missing data when estimating optimal DTRs. For survival outcomes,  methods for optimal DTRs estimation were developed by addressing missingness from right-censoring \citep{Simoneau_et_al_2020}.  For handling missing data in covariates, treatments and outcomes in general, \cite{Shortreed_et_al_2014} proposed a time-ordered nested conditional imputation strategy when estimating optimal DTRs in sequential multiple assignment randomized trials. \cite{Dong_et_al_2020} applied the augmented inverse probability weighting method to Q-learning and a generalized version of outcome-weighted learning when dealing with missing data due to patients' dropout. In simpler settings of estimating single-stage optimal treatment regimes, 
\cite{Shen_Hubbard_Linn_2023} further discussed the multiple imputation method proposed by \cite{Shortreed_et_al_2014}  for estimating and evaluating optimal individualized treatment rules that were not directly observed in the design. Also for the single-stage setting with missing data,  \cite{Huang_Zhou_2020} investigated the performance of an augmented inverse probability weighted estimator in the direct-optimization framework.  However, these existing methods all focus on ignorable missingness and could potentially yield sub-optimal DTRs when missing data for covariates, treatments or outcomes are nonignorable  \citep{Little_Rubin_2014}.  


In this research, we aim to address the issue of nonignorable missing covariates in the estimation of optimal DTRs, which is prominent in settings where patients are informatively monitored, e.g., in electronic medical records (EMR) databases or disease clinics. Our research was motivated by a recent investigation on optimal fluid strategy for sepsis patients in intensive care units (ICUs) \citep{Speth_et_al_2022}. 
Using EMR data from the Medical Information Mart for Intensive Care (MIMIC-III) \citep{Johnson_et_al_2016}, \cite{Speth_et_al_2022} estimated an optimal two-stage DTR for treating sepsis patients with fluid resuscitation in ICUs by incorporating baseline covariates such as demographics, time-varying covariates such as mechanical ventilation usage and vasopressor administration,  and the outcome measure of multi-organ failure at the first stage. However, according to the recommendations of the `Surviving Sepsis Campaign', hemodynamic variables should also be taken into account when reassessing the need for additional fluid administration following the initial resuscitation \citep{Rhodes_et_al_2017}. Unfortunately,  within the MIMIC-III database, certain hemodynamic variables, including blood pressure, respiratory rate, and body temperature, are subject to missingness at the conclusion of the initial resuscitation. Furthermore,  as highlighted by \cite{Awad_et_al_2017},  patients in the MIMIC-III database were informatively monitored since patients with more severe conditions, who typically presented with more abnormal physiological indicators than others, received more intensive monitoring.  Consequently, their measurements of hemodynamic variables were less likely to be missing. In this case, the missingness probabilities of the hemodynamic variables were directly related to their own values, indicating that the missing data mechanism was nonignorable missing. This characteristic could potentially lead to sub-optimal estimation of DTRs if existing methods are employed.

The identification and estimation of optimal DTRs become much more challenging when covariates are nonignorable missing. In the literature on handling  nonignorable missing covariates, it has been concluded that the parameters of an outcome regression model can be consistently estimated using complete case analysis as long as the missingness of the covariates is not directly related to the outcome \citep{Yang_Wang_Ding_2019}. Therefore, considering that Q-learning relies on consistently estimated outcome models and is both flexible and straightforward to implement, it offers a natural approach for accommodating nonignorable missing covariates when estimating optimal DTRs. 
However, when implementing the step-wise outcome regression algorithm in  Q-learning,  challenges arise not only from the missingness of the covariates but also from the missingness of the \emph{pseudo-outcomes}, even if  outcome variables are fully observed. This is because the Q-learning algorithm computes the pseudo-outcome in the preceding stage using covariates from the subsequent stage. As a result, if there are nonignorable missing covariates in the later stages, it can lead to nonignorable missing pseudo-outcomes in the earlier stages.  We provide a detailed illustration of this unique missing data problem with DTRs caused by backward-induction-induced nonignorable missing pseudo-outcomes in Section 2.3. 

It is well known that with nonignorable missing outcomes, outcome regression models are generally not identifiable; see examples provided by \cite{Miao_Ding_Geng_2016}. Inspired by the nonignorable missing outcome literature, we propose two weighted Q-learning approaches to tackling nonignorable missing  pseudo-outcomes in optimal DTR estimation. The first approach is to employ estimating equations to estimate a working model for missingness probability of the pseudo-outcome with valid nonresponse instrumental variables  \citep{Shao_Wang_2016}.  The pseudo-outcome estimates based on observed data are then weighted by the inverse of the estimated missingness probabilities in the Q-learning algorithm. In scenarios where nonresponse instrumental variables are not available, we specify a sensitivity parameter to quantify the residual association between the missing pseudo-outcome at stage $t$ and its conditional missingness probability given observed information up to stage $t$. Thus the missingness probabilities of the pseudo-outcomes are functions of the sensitivity parameter. We develop a practical approach to calibrating the sensitivity parameter  and then perform a sensitivity analysis to evaluate the impact of the pre-specified sensitivity parameter on the estimated optimal DTRs.
%\citep{Kim_Yu_2011}.

\section{Method}

\subsection{Setting and notation}

We consider an observational cohort with $n$ patients and finite $T$ treatment stages. Let subscripts $i=1,\ldots,n$, $t=1,\ldots,T$ denote patients and stages, respectively. Unless specified otherwise, we use capital letters to denote random variables and lowercase letters to indicate specific realizations of random variables.  Complete data for the patients are assumed independent and identically distributed and thus we suppress the patient-specific subscript $i$. Let $A_t\in \{-1, 1\}$ be the assigned treatment at stage $t$ for $t = 1, \ldots, T$.   $\bm{X}_t \in \mathbb{R}^{p_t}$ is a  $1 \times p_t$ vector of covariates measured before $A_t$, and $Y_t$ denotes the longitudinal %proximal
outcome measured after $A_t$ at stage $t$. 

The final outcome of interest is a pre-specified function (e.g., sum or maximum) of the longitudinal outcome in $T$ stages, $Y = f(Y_1,\ldots, Y_T)$. Higher values of $Y$ are assumed to be better. We define $\bm{H}_1=\bm{X}_1$, and $\bm{H}_t=\left(\bm{H}_{t-1}, A_{t-1}, Y_{t-1}, \bm{X}_t\right)$ for $t=2, \ldots, T$. Thus, $\bm{H}_t$ represents the information available before making the treatment decision at stage $t$. Let $\mathcal{H}_t$ denote the support of $\bm{H}_t$. A DTR $\bm{d}$ consists of a set of decision rules $(d_1, \ldots, d_T)$, where $d_t: \mathcal{H}_t \to \{-1, 1\}$ is a function that takes the observed history $\bm{h}_t$ as input and outputs a treatment decision at stage $t$. An optimal DTR is the set of decision rules $\left(d^{\mathrm{opt}}_1, \ldots, d^{\mathrm{opt}}_T\right)$ that maximizes the expectation of the final outcome. We formalize this definition using the potential outcome framework.

For $t=1,\ldots,T$, let $Y_t^*(\overline{\bm{a}}_t)$  denote the potential longitudinal outcome at stage $t$ and $\bm{X}_{t+1}^*(\overline{\bm{a}}_t)$  denote the potential covariates at stage $t+1$  if a patient, possibly contrary to fact, had received the treatment sequence $\overline{\bm{a}}_t = (a_1,\ldots,a_t)$  by stage $t$. We then define the set of potential outcomes under $\overline{\bm{a}}_t$ as $\bm{O}_t^*(\overline{\bm{a}}_t) = \{Y_1^*(a_1), \bm{X}_2^*(a_1),\ldots,Y_{t}^*(\overline{\bm{a}}_{t}),\bm{X}_{t+1}^*(\overline{\bm{a}}_{t}) \}$  for $t=1,\ldots,T-1$. The potential final outcome under a regime $\bm{d}$ is
\begin{align*}
Y^*(\bm{d})=\sum_{\overline{\bm{a}}_T} Y^*\left(\overline{\bm{a}}_T\right) \prod_{t=1}^T \mathbb I (d_t\left[\{\bm{X}_1,a_1,Y_1^*(a_1),\ldots,a_{t-1},Y_{t-1}^*(\overline{\bm{a}}_{t-1}),\bm{X}_{t}^*(\overline{\bm{a}}_{t-1})\}\right]=a_t),
\end{align*}
where $Y^*\left(\overline{\bm{a}}_T\right)=f\left\{Y_1^*\left(a_1\right), Y_2^*\left(\overline{\bm{a}}_2\right), \ldots, Y_T^*\left(\overline{\bm{a}}_T\right)\right\}$. 
Define the value of a regime $\bm{d}$ to be $V(\bm{d}) = E\{Y^*(\bm{d})\}$. Let $\mathcal{D}$ denote all feasible regimes. An optimal DTR, $\bm{d}^{\mathrm{opt}}\in \mathcal{D}$, satisfies that $V(\bm{d}^{\mathrm{opt}}) \geq V(\bm{d})$ for all $\bm{d} \in \mathcal{D}$.



\subsection{Q-learning with complete data}

In this section, we briefly review the Q-learning algorithm for estimating optimal DTRs with complete observational data. To estimate the value of candidate regimes from observed data, it is necessary to express the value of regimes solely in terms of observables rather than potential outcomes. This becomes feasible under the following, now standard, causal inference assumptions:
\begin{description}
    \item {\em Assumption} 1 (Consistency). $Y_t^*(\overline{\bm{a}}_t) = Y_t$ for $t = 1,\ldots,T$ and $\bm{X}_{t+1}^*(\overline{\bm{a}}_t) = \bm{X}_{t+1}$ for $t = 1,\ldots,T-1$ when $\overline{\bm{a}}_t$ are actually received.
    \item {\em Assumption} 2 (Sequential ignorability). $\{\bm{O}_{T-1}^*(\overline{\bm{a}}_{T-1}), Y_T^*(\overline{\bm{a}}_T) : \overline{\bm{a}}_T \in \bigotimes_{t=1}^{T} \{-1,1\}\} \upmodels A_t \mid \bm{H}_t$ for $t=1,\ldots,T$, where $\bigotimes$ denotes the Cartesian product.
    % This assumption denotes that all potential outcomes (Y_1^*(a_1), \bm{X}_2^*(a_1),\ldots,Y_{T-1}^*(\overline{\bm{a}}_{T-1}),\bm{X}_{T}^*(\overline{\bm{a}}_{t-1}),Y_T^*(\overline{\bm{a}}_T) ) are conditionally independent of $A_t$ given $\bm{H}_t$
    \item {\em Assumption} 3 (Positivity). $P(A_t=a_t\mid \bm{h}_t) > \epsilon >0$ for $a_t \in \{-1,1\}$, $\bm{h}_t \in  \mathcal{H}_t$, and $t=1,\ldots,T$, where $\epsilon$ is a positive constant. 
\end{description}

Under Assumptions 1-3, the value of a DTR can be identified from the observed data. Specifically, in the 2-stage scenario we have $E\{Y^*(\bm{d})\} = E(E[\{Y\mid \bm{H}_2,A_2 = d_2(\bm{H}_2)\}\mid \bm{H}_1,A_1 = d_1(\bm{H}_1)] )$.
% \begin{align*}
%     E[Y^*(\bm{d})] = E[E\{E(Y\mid \bm{H}_2,A_2 = d_2(\bm{H}_2))\mid \bm{H}_1,A_1 = d_1(\bm{H}_1)\} ] 
% \end{align*}
This repeated expectation form suggests that the optimal DTR can be computed by a backward-induction procedure, which is based on the following recursively defined Q-functions. % in canonical Q-learning.
% :
% \begin{align*}
%  Q_T\left(\bm{h}_T, a_T\right) &= E\left(Y \mid \bm{H}_T=\bm{h}_T, A_T=a_T\right), \\
%  Q_t\left(\bm{h}_t, a_t\right) &= E\left\{\max _{a_{t+1} \in \{-1,1\}} Q_{t+1}\left(\bm{H}_{t+1}, a_{t+1}\right) \mid \bm{H}_t=\bm{h}_t, A_t=a_t\right\}, 
% \end{align*}
% for $ t=T-1,\ldots,1$. 
At the final stage $T$, $Q_T\left(\bm{h}_T, a_T\right) = E\left(Y \mid \bm{H}_T=\bm{h}_T, A_T=a_T\right)$. For $ t=T-1,\ldots,1$, $Q_t\left(\bm{h}_t, a_t\right) = E\{\max_{a_{t+1} \in \{-1,1\}} Q_{t+1}(\bm{H}_{t+1}, a_{t+1}) \mid \bm{H}_t=\bm{h}_t, A_t=a_t\}.$

The true Q-functions are unknown and must be estimated from the data. Since Q-functions represent conditional expectations, it is natural to estimate them using regression models. Note that in these models, the response variable is $\max _{a_{t+1} \in \{-1,1\}} Q_{t+1}\left(\bm{H}_{t+1}, a_{t+1}\right)$ instead of $Y_t$ ($t = 1,\ldots,T-1$). These response variables are commonly referred to as pseudo-outcomes, and we define $Y_{pse,T} = Y$ and $Y_{pse,t} = \max _{a_{t+1} \in \{-1,1\}} Q_{t+1}\left(\bm{H}_{t+1}, a_{t+1}\right)$ for $t = 1,\ldots,T-1$. Subsequently, $Q_t(\bm{h}_t,a_t) = E\left( Y_{pse,t}\mid \bm{h}_t,a_t \right)$ for $t = 1,\ldots,T$. Since $a_t \in \{-1,1\}$, any specified model for $Q_t(\bm{h}_t,a_t)$ can be decomposed as $Q_t(\bm{h}_t,a_t;\bm{\theta}_t) = q_{t,0}(\bm{h}_t; \bm{\beta}_t) + a_t q_{t,1}(\bm{h}_t;\bm{\psi}_t)$, where $\bm{\theta}_t = (\bm{\beta}_t, \bm{\psi}_t)$, $q_{t,0}(\bm{h}_t) = \frac{Q_t(\bm{h}_t,1) + Q_t(\bm{h}_t,-1)}{2}$, and $q_{t,1}(\bm{h}_t)  = \frac{Q_t(\bm{h}_t,1) - Q_t(\bm{h}_t,-1)}{2}$ for $t=1\ldots,T$. As a result, the Q-function is divided into two components: (1) a stage $t$ treatment-free component $ q_{t,0}(\bm{h}_t; \bm{\beta}_t)$ which depends on (a subset of) the history $\bm{h}_t$  before stage $t$ but not on the stage $t$ treatment $a_t$, and (2) a stage $t$ treatment effect component $a_tq_{t,1}(\bm{h}_t;\bm{\psi}_t)$ which depends on (a potentially different subset of) $\bm{h}_t$ and specifically includes the main effect of treatment $a_t$ and its interactions with tailoring variables (i.e., effect modifiers). The blip function, as defined in \cite{Chakraborty_Moodie_2013}, equals to $a_t q_{t,1}(\bm{h}_t;\bm{\psi}_t) + q_{t,1}(\bm{h}_t;\bm{\psi}_t)$ under this setting. Therefore, we refer to $\bm{\psi}_t$ as the blip function parameters. Using this modeling approach, we have $d_t^{\mathrm{opt}}(\bm{h}_t) = 2\mathbb I \left\{q_{t,1}(\bm{h}_t;\bm{\psi}_t)>0\right\} -1$. 
Under Assumptions 1-3, the estimator $\widehat{\bm{\theta}}_t$ of $\bm{\theta}_t$ can be recursively constructed by regressing $Y_{pse,t}$ on $\bm{H}_t$ and $A_t$ . Then, the Q-learning estimator of $d^{\mathrm{opt}}_t$ can be obtained by $\widehat{d}_t^{\mathrm{opt}}\left(\bm{h}_t\right) = 2\mathbb{I}\left\{q_{t,1}(\bm{h}_t;\widehat{\bm{\psi}}_t)>0\right\} -1$ for $t=1,\ldots,T$.


\subsection{Q-learning for DTRs with  nonignorable missing covariates: why complete case analysis fails at earlier stages while working at the final stage?}

We focus on the scenario that the treatment and longitudinal outcome, $(A_t, Y_t)$, are fully observed while the covariates, $\bm{X}_t$, contain missing values. This scenario is frequently encountered in practice because treatments are typically well-documented and the longitudinal outcomes of interest are usually key clinical disease activity indices which are regularly measured. Moreover, the probability of a covariate being missing is likely to depend on the actual value of the covariate itself due to informative monitoring of patients. 
For instance, in the MIMIC-III data, patients who displayed symptoms of abnormal heart rates such as chest pain were more likely to have their heart rates recorded. Therefore, the absence of heart rate measurements could indicate nonignorable missingness, since patients with normal heart rates were less likely to be frequently monitored.

Let $R_t$ represent the missingness indicator for $\bm{X}_t$ such that $R_{t} = 1$ if all variables in $\bm{X}_t$ are fully observed and $R_{t} = 0$ if at least one component of $\bm{X}_t$ is missing. Define $\overline{\bm{R}}_t = (R_1,\ldots, R_t)$ as the history of missingness indicators up to stage $t$. We use $\bm{1}_{t}$ to denote a $1 \times t$ vector in which all elements are equal to $1$ for $t=1, \ldots, T$. In ICU settings, the effect of a treatment is commonly evaluated by the patient's longitudinal outcomes (e.g., physiological symptoms) observed several hours after treatment administration, while the treatment decision for a patient is based on the patient's medical history. Because of the temporal ordering of these covariates, it is plausible that the vector of missingness indicators $\overline{\bm{R}_t}$ does not directly depend on the patient's variables measured after receiving the current stage treatment $A_t$. Thus we make the following assumption:
\begin{description}
    \item {\em Assumption} 4 (Future-independent missingness). $\overline{\bm{R}_t} \upmodels (Y_t, \bm{X}_{t+1}, R_{t+1},  \ldots, \bm{X}_T,R_T,A_T,Y_T) \mid (\bm{H}_t, A_t)$  for $t=1,\ldots,T$.\label{A4}
\end{description}
Figure \ref{DAG} is a directed acyclic graph  illustrating Assumptions 2 and 4 in a 2-stage scenario. In this example, $\bm{H}_1 = \bm{X}_1$, $\bm{H}_2 = (\bm{X}_1,A_1,Y_1,\bm{X}_2)$. $R_1$ and $R_2$ are missing indicators for $\bm{X}_1$ and $\bm{X}_2$, respectively. Graphically, $A_1$ has no common parents with ($Y_1, \bm{X}_2, A_2, Y_2$) except for $\bm{H}_1$, and $A_2$ has no common parents with $Y_2$ except for $\bm{H}_2$, encoding Assumption 2. Besides, $R_1$ is d-separated from $(Y_1, \bm{X}_2,R_2, A_2, Y_2)$ by $(\bm{H}_1,A_1)$, and $(R_1, R_2)$ is d-separated from $Y_2$ by $(\bm{H_2},A_2)$ \citep{VanderWeele_Robins_2007}, encoding Assumption 4. 
\begin{figure}[htbp]
    \centering
    \includegraphics[scale=0.45]{DAG_1113.png}
    \caption{A direct acyclic graph illustrating Assumptions 2 and 4}
    \label{DAG}
\end{figure}

In the single-stage scenario, Assumption 4 degenerates to the outcome-independent missingness assumption \citep{Yang_Wang_Ding_2019}, which posits that the missingness probability of covariates or confounders is conditionally independent of the outcome, given all variables except for the outcome, i.e., $R_1 \upmodels Y_1\mid (\bm{X}_1, A_1)$. Under this assumption, we have $P(Y_1\mid \bm{X}_1, A_1) = P(Y_1\mid \bm{X}_1, A_1, R_1=1)$, so that we can obtain unbiased parameter estimates of the outcome model with complete case analysis. Consequently, in the single-stage scenario, Q-learning based on complete case analysis yields consistently optimal individualized treatment rules. In this research, considering the plausibility of  conditional independence ensured by temporal ordering of the variables, we adapt the outcome-independent missingness assumption to longitudinal settings with multiple stages and make Assumption 4. 

Given $\bm{H}_t$ and $A_t$, the pseudo-outcome $Y_{pse,t}$ only depends on $Y_t$ and the variables measured after stage $t$. Therefore, under Assumption 4, $Y_{pse,t} \upmodels \overline{\bm{R}}_t \mid (\bm{H}_t, A_t$) and we have $P(Y_{pse,t} \mid \bm{H}_t, A_t) = P(Y_{pse,t} \mid \bm{H}_t, A_t, \overline{\bm{R}}_t = \bm{1}_t)$ for $t = 1,\ldots,T$.
Thus, we can obtain consistent estimates of $Q_t\left(\bm{h}_t, a_t\right)$ by utilizing data from patients with complete data up to stage $t$ if there are no missing values in  $Y_{pse,t}$. For example, at stage $T$, $Y_{pse,T}=Y$ is fully observed. Thus we can use complete case analysis to estimate $Q_T(\bm{h}_T, a_T)$.

However, conditional on $\overline{\bm{R}}_t=\bm{1}_t$ (i.e., for complete cases up to stage $t$), although there is no missingness in $\bm{H}_t$, the pseudo-outcome $Y_{pse,t}$ may contain missing values because $Y_{pse,t}$  depends on $\bm{X}_{t+1}$, which can be missing  (i.e.,  $R_{t+1} =0$) when $\overline{\bm{R}}_t=\bm{1}_t$. We illustrate this backward-induction-induced missing pseudo-outcome problem in Figure \ref{pseudo-outcome missing pattern} with a hypothetical example coming from the same setting of our simulation study in Section 3.  In this example, all variables for Patient 3 at stage one are fully observed. Nevertheless, $Y_{pse,1}$ for Patient 3 is unavailable because $X_{22}$, a covariate required for computing $Y_{pse,1}$, is missing. 

\begin{figure}[htbp]
    \centering
    \includegraphics[scale=0.27]{missing_indicator_4.png}
    \caption{A hypothetical example illustrating the relationship between the missing data pattern of covariates and that of pseudo-outcomes. Gray areas represent missing information. The setting of the variables is the same as that in Simulation 1 in Web Appendix H.}
    \label{pseudo-outcome missing pattern}
\end{figure}

Let $R_{pse,t}$ denote the missingness indicator of $Y_{pse,t}$ for $t=1,\ldots,T$. Because the missingness probabilities of $\bm{X}_{t+1}$ depend on their own values, and $Y_{pse,t}$ is computed based on $\bm{X}_{t+1}$, it is likely that $Y_{pse,t}$ is nonignorably missing. This backward-induction-induced nonignorable missing pseudo-outcome problem would make identification and consistent estimation of Q-functions challenging. A complete case analysis for Q-learning would likely result in suboptimal DTR. In Section 3, we will use simulations to demonstrate this phenomenon. 

\subsection{Weighted Q-learning for DTRs with nonignorable missing pseudo-outcomes: a nonresponse instrument variable approach}

We propose two approaches for inverse probability weighted Q-learning to tackle the nonignorable missing pseudo-outcome problem. The first involves utilizing valid nonresponse instrumental variables. Specifically, we consider a semiparametric model for $R_{pse,t}$:
\begin{description}
    \item {\em Assumption} 5 (semiparametric working model for missingness probability of the pseudo-outcomes). For $t = 1,\ldots,T-1$, 
\begin{align}\label{missingnessmodel}
    P(R_{pse,t}=1\mid \bm{h}_t,a_t,y_{pse,t},\overline{\bm{R}}_t=\bm{1}_t) = \pi(\bm{u}_t,y_{pse,t};\gamma_t)= \frac{1}{1+\exp\{s_t(\bm{u}_t)+\gamma_t y_{pse,t}\}},
\end{align}
\end{description}where $\bm{U}_t \subset (\bm{H}_t,A_t)$ is a proper subset of $(\bm{H}_t,A_t)$, $P(Y_{pse,t}\mid \bm{H}_t, A_t) \neq P(Y_{pse,t}\mid \bm{U}_t)$, and $s_t(\cdot)$ is an unknown and unspecified function of $\bm{u}_t$. %For convenience, we use $\pi(\bm{u}_t,y_{pse,t};\gamma_t)$ to denote $\frac{1}{1+\exp\{s_t(\bm{u}_t)+\gamma_t y_{pse,t}\}}$.

This semiparametric model assumption does not impose any parametric constraint on the function of $\bm{u}_t$ in working model \eqref{missingnessmodel}. Instead, it restricts the parametric form of $y_{pse,t}$ and requires that there is no interaction between $\bm{u}_t$ and $y_{pse,t}$ in working model \eqref{missingnessmodel}. 
It is noteworthy that the variables observed after stage $t$ may exhibit an association with $R_{pse,t}$ even after conditioning on ($\bm{h}_t,a_t,y_{pse,t},\overline{\bm{R}}_t=\bm{1}_t$) and therefore model \eqref{missingnessmodel} may not be the true data-generating model for $R_{pse,t}$. However,  weighting by the inverse of  $P(R_{pse,t}=1\mid \bm{h}_t,a_t,y_{pse,t},\overline{\bm{R}}_t=\bm{1}_t)$ would suffice to balance the distribution of the variables between the patients with and without missing $y_{pse,t}$ for the purpose of estimating $Q_t(\bm{h}_t,a_t)$ and thus lead to consistent estimation of the Q-functions after weighting. 

Assumption 5 also implies that $(\bm{H}_t,A_t)$ can be decomposed as $(\bm{H}_t,A_t) = (\bm{U}_t,\bm{Z}_t)$ such that $\bm{Z}_t \nupmodels Y_{pse,t} \mid \bm{U}_t$ and $\bm{Z}_t \upmodels R_{pse,t} \mid  \bm{U}_t, Y_{pse,t}$ when $\overline{\bm{R}}_t=\bm{1}_t$. That is, conditional on $\bm{U}_t$, $\bm{Z}_t$ is a predictor of the pseudo-outcome $Y_{pse,t}$, but given $\bm{U}_t$ and $Y_{pse,t}$, $\bm{Z}_t$ is conditionally independent of the missingness indicator $R_{pse,t}$.
$\bm{Z}_t$ is referred to as the nonresponse instrumental variables at stage $t$. The nonresponse instrumental variables, sometimes known as the `shadow variables', are commonly employed in the literature for nonignorable missing outcomes.  It has been demonstrated in previous studies of nonignorable missing outcomes \citep{Shao_Wang_2016, Miao_Tchetgen_2016} that  nonresponse instrumental variables facilitate the identification of the working model \eqref{missingnessmodel} and the outcome model. In certain scenarios, it may not be difficult to find valid nonresponse instrumental variables. For example, gender and weight are associated with the severity of sepsis in ICU patients \citep{kalani2020sepsis, sakr2013influence}. It is plausible that they are independent of the missing data mechanism of hemodynamic variables, given other pre-treatment covariates of the patients  (e.g., comorbidities) and the severity of sepsis represented by $Y_{pse,t}$. These characteristics render gender and weight suitable candidates for nonresponse instrumental variables.


If Assumptions 5 holds, to estimate the working  model~\eqref{missingnessmodel}, we have to overcome the challenge that $s_t(\cdot)$ is nonparametric. Following \cite{Shao_Wang_2016}, we consider the idea of profiling and obtain the following kernel estimate of $s_t(\cdot)$ for a given $\gamma_t$.
\begin{align}
    \exp\{\widehat{s}_{\gamma_t,t}(\bm{u}_t) \} =  \frac{\sum_{i=1}^{n} \left(1-r_{pse,t,i}\right)  K_{c_t}\left(\bm{u}_t-\bm{u}_{t,i}\right) \mathbb{I}(\overline{\bm{R}}_{t,i}=\bm{1}_t) }{\sum_{i=1}^{n}  r_{pse,t,i}   \exp \left(\gamma_t y_{pse,t,i}\right) K_{c_t}\left(\bm{u}_t-\bm{u}_{t,i}\right) \mathbb{I} (\overline{\bm{R}}_{t,i}=\bm{1}_t)}, \label{2.1}
\end{align}
where $K_{c_t}(\cdot) = c_t^{-1}K(\cdot/c_t)$, with $K(\cdot)$ being a symmetric kernel function and $c_t$ a bandwidth. 
For a fixed $\gamma_t$, the working model \eqref{missingnessmodel} can be expressed as $\widehat{\pi}_t(\bm{u}_t,y_{pse,t};\gamma_t) = [1+\exp\{\widehat{s}_{\gamma_t,t}(\bm{u}_t) + \gamma_t y_{pse,t}\} ]^{-1}$.  Then, we can estimate $\gamma_t$ with the following estimating equation,
\begin{align}
    \widehat{E} \left[ \mathbb{I} (\overline{\bm{R}}_t=\bm{1}_t) \bm{l}_t(\bm{z}_t)\left\{\frac{r_{pse,t}}{\widehat{\pi}_t(\bm{u}_t,y_{pse,t}; \gamma_t)} -1 \right\} \right]  = \bm{0}, \label{2.2} 
\end{align}
where $\widehat{E}[\cdot]$ denotes the empirical mean, $\bm{l}_t(\bm{z}_t) = (l_{t,1}(\bm{z}_t), \ldots, l_{t, L}(\bm{z}_t))$ is a user-specified differentiable vector function of the nonresponse instrumental variables of dimension $L$, $L\geq2$. 

Note that in estimating equations \eqref{2.2}, if $\bm{z}_t$ are constants, then the above equations are not solvable because they are under-identified. Nevertheless, when $L > 2$ and the equations are over-identified, we can employ the two-step generalized method of moments (GMM) to estimate $\gamma_t$. Let $b_{t,j}(r_{pse,t}, \bm{u}_t, y_{pse,t};\gamma_t) = \mathbb{I} (\overline{\bm{R}}_t=\bm{1}_t) l_{t,j}(\bm{z}_t)\left\{\frac{r_{pse,t}}{\widehat{\pi}_t(\bm{u}_t,y_{pse,t}; \gamma_t)} -1 \right\}$ for $j=1,\ldots,L$ and $\bm{B}_t(\gamma_t) = (\widehat{E}[b_{t,1}],\widehat{E}[b_{t,2}],\ldots,\widehat{E}[b_{t,L}] )$. Then, the first-step GMM estimator of $\gamma_t$ is $\widehat{\gamma}^{(1)}_t=\underset{\gamma_t}{\arg\min } \; \bm{B}_t(\gamma_t) \bm{B}_t(\gamma_t)^{\mathrm{T}}$.


Based on the first-step GMM estimator, we can build a more efficient estimator of $\gamma_t$ with a weight matrix $\bm{W}_t$. Let $\bm{W}_{t}$ be the inverse matrix of the $L \times L$ matrix with the $\left(j, j^{\prime}\right)$ element as $\frac{1}{n} \sum_{i=1}^{n} b_{t,j}(r_{pse,t}, \bm{u}_t, y_{pse,t};\widehat{\gamma}^{(1)}_t) b_{t,j'}(r_{pse,t}, \bm{u}_t, y_{pse,t};\widehat{\gamma}^{(1)}_t)$.
Then the second-step GMM estimator of $\gamma_t$ is $\widehat{\gamma}_t = \underset{\gamma_t}{\arg\min }  \; \bm{B}_t(\gamma_t) \bm{W}_{t} \bm{B}_t(\gamma_t)^{\mathrm{T}}$.

With the estimated working model for the missingness probability, we can employ the weighted least square method to obtain consistent estimates of the Q-function parameters. Subsequently, the optimal treatment at stage $t$ can be determined based on the estimated parameters. We henceforth refer to this approach as the estimating equation (EE) method.


\subsection{Weighted Q-learning for DTRs with nonignorable missing pseudo-outcomes: a sensitivity analysis approach}

In real data applications, it might be challenging to find valid nonresponse instrument variables for every stage based on domain knowledge. 
In such scenarios, we develop a sensitivity analysis approach 
%\tb{to assess the robustness of optimal DTR estimation against the departures from ignorable missingness}. 
for the conditional missingness probability of $Y_{pse,t}$. 
Specifically, we slightly modify Assumption 5 and introduce the following assumption.
%\tb{it can be challenging to find nonresponse instrument variables for every stage based on domain knowledge and  their validity is unverifiable from observed data.}
% \tb{In practice, valid nonresponse instrument variables may not be available or justified based on domain knowledge.}
% %In such scenarios, we develop a sensitivity analysis approach for the conditional missingness probability of $Y_{pse,t}$. 
% \tb{%Therefore, sensitivity analysis is required.   %serves as a valuable complement for the EE method proposed above,  enhancing the reliability of the estimated optimal DTRs. Therefore, 
% In this case, we develop a sensitivity analysis approach %for the conditional missingness probability of $Y_{pse,t}$ 
% %to assessing the impact of departure from ignorable missingness on estimated DTRs.
% to assessing the impact of departure from key missingness assumptions and providing a plausible set of estimated DTRs.
% }
%Specifically, we slightly modify Assumption 5 and introduce the following assumption.
%\citep{european2010guideline}
\begin{description}
    \item {\em Assumption} 6 (Semiparametric working model for missingness probability  without nonresponse instrumental variables).
    \begin{align} \label{missingmodel2}
        P(R_{pse,t}=1\mid \bm{h}_t, a_t, y_{pse,t},\overline{\bm{R}}_t = \bm{1}_t) =\pi(\bm{h}_t,a_t,y_{pse,t};\gamma_t')= \frac{1}{1+\exp\{s'_t(\bm{h}_t,a_t)+\gamma'_t y_{pse,t}\}},
    \end{align}
    where $s'_t(\cdot)$ is an unknown and unspecified function of $(\bm{h}_t,a_t)$ and $\gamma'_t$ is  the pre-specified sensitivity parameter. Note that here we include all historical variables in~\eqref{missingmodel2}. %For brevity, we use $\pi(\bm{h}_t,a_t,y_{pse,t};\gamma_t')$ to denote $\frac{1}{1+\exp\{s_t(\bm{h}_t,a_t,y_{pse,t})+\gamma_t' y_{pse,t}\}}$
\end{description}

%The sensitivity parameter $\gamma'_t$ quantifies the residual effect of $Y_{pse,t}$ on $R_{pse,t}$ after adjusting for $\bm{h}_t$ and $a_t$ for patients with complete records up to stage $t$. The working model \eqref{missingmodel2} can be identified with a fixed $\gamma'_t$ \citep{Kim_Yu_2011}. We can obtain the kernel estimate of $\exp\{s'_{\gamma'_t,t}(\bm{h}_t)\}$ in a similar manner as in equation \eqref{2.1}.

% \begin{align}
%     \exp\{\hat{s}'_{\gamma'_t,t}(\bm{h}_t) \} =  \frac{\sum_{i=1}^{n} \left(1-r_{pse,t,i}\right)  K'_{c'_t}\left(\bm{h}_t-\bm{h}_{t,i}\right) \mathbb{I}(\overline{\bm{R}}_{t,i}=\bm{1}_t) }{\sum_{i=1}^{n}  r_{pse,t,i}   \exp \left(\gamma'_t y_{pse,t,i}\right) K'_{c'_t}\left(\bm{h}_t-\bm{h}_{t,i}\right) \mathbb{I} (\overline{\bm{R}}_{t,i}=\bm{1}_t)}, \label{2.3} %\tag{2.3}
% \end{align}


The sensitivity parameter $\gamma'_t$ quantifies the residual effect of $Y_{pse,t}$ on $R_{pse,t}$ after adjusting for $\bm{h}_t$ and $a_t$ for patients with complete covariates up to stage $t$. \cite{Kim_Yu_2011} proposed a similar semiparametric exponential tilting model for nonignorable missing outcomes in cross-sectional settings, where external data were used to estimate their sensitivity parameter. In more general scenarios where external data are unavailable, the specification of sensitivity parameters in  nonignorable missingness models (and  in related problems of sensitivity analysis for unmeasured confounding) remains an active  area of research. In Web Appendix A, we first discuss how to interpret the sign of $\gamma'_t$ using domain knowledge, and then develop a simulation-based approach to calibrating the magnitude of  $\gamma'_t$, building upon the methods by \cite{yin2019simulation}.

After $\gamma_t'$ is specified, we can obtain the kernel estimate of $\exp\{s'_{\gamma'_t,t}(\bm{h}_t,a_t)\}$ in a similar manner as in equation \eqref{2.1}:
%Let $\bm{G}_t = (\bm{H}_t,A_t)$, we have}
\begin{align}
    \exp\{\widehat{s}'_{\gamma'_t,t}(\bm{h}_t,a_t) \} =  \frac{\sum_{i=1}^{n} \left(1-r_{pse,t,i}\right)  K'_{c'_t}\left(\bm{g}_t-\bm{g}_{t,i}\right) \mathbb{I}(\overline{\bm{R}}_{t,i}=\bm{1}_t) }{\sum_{i=1}^{n}  r_{pse,t,i}   \exp \left(\gamma'_t y_{pse,t,i}\right) K'_{c'_t}\left(\bm{g}_t-\bm{g}_{t,i}\right) \mathbb{I} (\overline{\bm{R}}_{t,i}=\bm{1}_t)}, \label{2.3} %\tag{2.3}
\end{align}
where $\bm{g}_t = (\bm{h}_t,a_t)$ and $K'_{c'_t}(\cdot) = {c'_t}^{-1}K'(\cdot/c'_t)$, with $K'(\cdot)$ being a symmetric kernel function and $c'_t$ a bandwidth.

Subsequently, we can employ inverse probability weighting  to obtain consistent estimates of the Q-functions and the optimal DTRs. If the working model~\eqref{missingmodel2} is correctly specified (including the sensitivity parameter $\gamma'_t$)  and certain regularity conditions for $\exp\{s'_{\gamma'_t,t}(\bm{h}_t)\}$ are satisfied, along with correctly specified Q-functions, the parameter estimators in the Q-functions are consistent. As the true value of  $\gamma'_t$ is unknown, we conduct a sensitivity analysis by varying the value of  $\gamma'_t$ in the determined range to assess its  impact  on the Q-learning results. %Given that the robustness of this method hinges on sensitivity analysis,
 We refer to this approach as the SA method.


\subsection{Summary of the estimation steps}\label{steps}
We summarize the steps of the proposed weighted Q-learning approach as follows: 


\begin{enumerate}
    \item  Compute $y_{pse,T} = f(y_{1},y_{2}, \ldots, y_{T})$. Let $\bm{M}_T(\bm{\theta}_T) =\widehat{E}\left[ \mathbb{I}(\overline{\bm{R}}_{T}=\bm{1}_T) \frac{\partial \left[ \left\{  
  Q_T(\bm{h}_T,a_T;\bm{\theta}_T) - y_{pse,T}\right\} ^2\right]}{\partial \bm{\theta}_T} \right]$. Then obtain $\widehat{\bm{\theta}}_T$ by solving the estimating equation $\bm{M}_T(\bm{\theta}_T) = \bm{0}$
    \item  For $t = T-1,\ldots, 1$, repeat the following steps recursively:
    \begin{enumerate}
        \item Obtain stage $t$ pseudo-outcome estimate as $\widehat{y}_{pse,t} = \max\limits_{a_{t+1}} \; Q_{t+1}(\bm{h}_{t+1},a_{t+1};\widehat{\bm{\theta}}_{t+1})  $, denote the
        missingness indicator of $\widehat{y}_{pse,t}$ as $\widehat{R}_{pse,t}$.
         %missingness indicator of $\widehat{Y}_{pse,t}$ as $\widehat{R}_{pse,t}$. 
        \item Estimate $P(\widehat{R}_{pse,t}=1\mid \bm{h}_t,a_t, \widehat{y}_{pse,t}, \overline{\bm{R}}_t=\bm{1}_t)$  using the EE or the SA method and  denote the estimated probability as $\widehat{\pi}_t$.% with corresponding working models for the missingness probability of the estimated  pseudo-outcome.
        \item Let $\bm{M}_t(\bm{\theta}_t) = \widehat{E} \left(\frac{\widehat{r}_{pse,t}}{\widehat{\pi}_{t}} \mathbb{I}(\overline{\bm{R}}_{t}=\bm{1}_t) \frac{\partial \left[ \left\{ Q_t(\bm{h}_t,a_t;\bm{\theta}_t) - \widehat{y}_{pse,t}  \right\}^2\right]}{\partial \bm{\theta}_t} \right)$. Solve the estimating equation \\$\bm{M}_t(\bm{\theta}_t) = \bm{0}$ to obtain $\widehat{\bm{\theta}}_t$.
    \end{enumerate}
    \item  Obtain the optimal  regime at stage $t$  by $\widehat{d}_t^{\mathrm{opt}}\left(\bm{h}_t\right) =  \argmax\limits_{a_t}Q_t(\bm{h}_t,a_t;\widehat{\bm{\theta}}_t)$ for $t=1,\ldots,T$. 
\end{enumerate}


{\em Remark 1.} Recall that, since $a_t\in\{-1,1\}$,  $Q_t(\bm{h}_t,a_t;\widehat{\bm{\theta}}_t)$ can be written as $ q_{t,0}(\bm{h}_t; \widehat{\bm{\beta}}_t) + a_t q_{t,1}(\bm{h}_t;\widehat{\bm{\psi}}_t)$. Subsequently,  $\widehat{y}_{pse,t}$ can be calculated by $q_{t+1,0}(\bm{h}_{t+1}; \widehat{\bm{\beta}}_{t+1}) + |q_{t+1,1}(\bm{h}_{t+1};\widehat{\bm{\psi}}_{t+1})|$ and $\widehat{d}_t^{\mathrm{opt}}\left(\bm{h}_t\right)$ is determined by the sign of $q_{t,1}(\bm{h}_t;\widehat{\bm{\psi}}_t)$.

{\em Remark 2.} The pseudo-outcome $Y_{pse,t}$ is determined by the specified parametric form of $Q_{t+1}(\cdot)$ and the true values of $\bm{\theta}_{t+1}$ for $t=1,\ldots,T-1$. Since the  true value of $\bm{\theta}_{t+1}$ is unknown in practice, we use $\widehat{Y}_{pse,t}$ calculated by the parameter estimates $\widehat{\bm{\theta}}_{t+1}$ in  step (2). In this scenario,  $\widehat{Y}_{pse,t}$ differs slightly from $Y_{pse,t}$, which leads $\widehat{\pi}_t$ to differ slightly from $\pi(\bm{u}_t,y_{pse,t})$ and $\pi(\bm{h}_t,a_t,y_{pse,t})$, when the EE and SA methods are used respectively. Nevertheless, if all the models are correctly specified, $\widehat{\bm{\theta}}_{t+1}$ serves as a consistent estimator for $\bm{\theta}_{t+1}$. Subsequently, $\widehat{Y}_{pse,t}$  converges to $Y_{pse,t}$. It follows that the estimated optimal DTR $\widehat{\bm{d}}$ will also converge to the true optimal DTR $\bm{d}$. We formalize these asymptotic properties in the following section. %In Section 3, we also conduct simulations to assess the robustness of the proposed EE method against mild model misspecification, e.g., when the chosen instrumental variable is weakly associated with  $R_{pse,t}$ after conditioning on $\bm{U}_t$ and $Y_{pse,t}$.

Remark 3.
$Y_{pse,t}=\max_{a_{t+1}\in \{-1,1\}} Q_{t+1}(\bm{H}_{t+1},a_{t+1})$ implies that 
%$R_{pse,t} \geq \mathbb{I}(\overline{\bm{R}}_{t+1} = \bm{1}_{t+1} )$. That is, 
%if a patient's pseudo-outcome is missing at stage $t$, then there is at least one missing member in $\bm{H}_{t+1}$. 
%the missingness of $Y_{pse,t}$ arises from the missingness of $\bm{H}_{t+1}$.
the missingness of $Y_{pse,t}$ is determined by the model for $Q_{t+1}(\cdot)$ and the missingness of $\bm{H}_{t+1}$.
If all members of $\bm{H}_{t+1}$ are employed in the calculation of $Y_{pse,t}$, then $R_{pse,t} = \mathbb{I}(\overline{\bm{R}}_{t+1} = \bm{1}_{t+1} )$. Therefore, the missingness of covariates at stage $t+2, \ldots, T$ does not directly affect $R_{pse,t}$. Nevertheless, if nonignorable missing covariates in the later stages are not handled properly, it not only directly impacts the estimation at those stages but also has an influence on earlier stages. This is because bias would propagate and accumulate during the backward-induction procedure through the pseudo-outcomes calculated with biased Q-functions. 
%Consequently, although the proportions of missing pseudo-outcomes at existing non-final stages would not be directly affected if more future stages are added, the overall correct classification rates of estimated optimal DTRs would become lower in this scenario.
%Therefore, the problem of nonignorable missing pseudo-outcomes could exacerbate as the total number of stages increases. We illustrate this in Simulation 4.
Therefore, in this scenario increasing the total number of stages could lead to worse performance of the estimated optimal DTRs. 
We investigate this issue further in Simulation 4  in Web Appendix H.


% \subsection{Asymptotic properties}

% The consistency and asymptotic normality of the estimators in Section~\ref{steps}  can be established under suitable regularity conditions for estimating equations and an additional assumption that the optimal treatment is unique for all subjects at all stages. Note that the parameters are estimated separately using estimating equations recursively at each stage. Therefore, it is natural to establish the asymptotic properties of the stage-specific parameters recursively as well. For simplicity, we will focus on the two-stage setting, but extensions to the general $T$-stage setting follow directly. 

% \begin{theorem}
% In the 2-stage scenario, if Assumptions 1-4 hold, the Q-functions at stage 2 are Lipschitz continuous and correctly specified, and $\bm{M}_2(\bm{\theta}_2) = \bm{0}$ has a unique solution, %then
% \begin{align*}
%     \widehat{\bm{\theta}}_2 - \bm{\theta}_{2}^* &\stackrel{d}{\longrightarrow} N(\bm{0}, \Sigma_{\bm{\theta}_2}) , \\
%     \widehat{d}_2^{\mathrm{opt}}(\bm{h}_2) & \stackrel{p}{\longrightarrow} d_2^{\mathrm{opt}}(\bm{h}_2) \text{ for all } \bm{h}_2 \in \mathcal{H}_2,\\
%     \widehat{Y}_{pse,1}  &\stackrel{p}{\longrightarrow} Y_{pse,1},
% \end{align*}
% where $\bm{\theta}_{2}^*$ stands or the true values of $\bm{\theta}_{2}$. $\widehat{\bm{\theta}}_{2}$, $\widehat{d}_{2}^{\mathrm{opt}}(\cdot)$, and $\widehat{Y}_{pse,1}$ denote the estimators of $\bm{\theta}_{2}$, $d_2^{\mathrm{opt}}(\cdot)$, and $Y_{pse,1}$, respectively. Notations $\stackrel{d}{\longrightarrow}$ and  $\stackrel{p}{\longrightarrow}$ stand for converge in distribution and converge in probability, respectively. The detailed form of $\Sigma_{\bm{\theta}_2}$ is given in Web Appendix  B.
% \end{theorem}



% \begin{theorem}
% In the two-stage scenario, if Assumptions 1-5 hold, all models are Lipschitz continuous and correctly specified, and the estimating equations for the Q-functions and the estimating equation \eqref{2.2} each have a unique solution, then under certain additional regularity conditions, 
% \begin{align*}
%     \widehat{\bm{\theta}}_{1,EE} - \bm{\theta}_{1}^* &\stackrel{d}{\longrightarrow} N(\bm{0}, \Sigma_{\bm{\theta}_1}) , \\
%     \widehat{d}_{1,EE}^{\mathrm{opt}}(\bm{h}_1) & \stackrel{p}{\longrightarrow} d_1^{\mathrm{opt}}(\bm{h}_1) \text{ for all } \bm{h}_1 \in \mathcal{H}_1,
% \end{align*}
% where $\bm{\theta}_{1}^*$ stands or the true values of $\bm{\theta}_{1}$. $\widehat{\bm{\theta}}_{1,EE}$ and $\widehat{d}_{1,EE}^{\mathrm{opt}}(\cdot)$ denote the estimators of $\bm{\theta}_{1}$ and $d_1^{\mathrm{opt}}(\cdot)$ with the working model \eqref{missingnessmodel} estimated by the EE method, respectively. The regularity conditions and the detailed form of $\Sigma_{\bm{\theta}_1}$ are given in Web Appendix  C.
% \end{theorem}


% \begin{theorem}
% In the two-stage scenario, if Assumptions 1-4 and 6 hold, all models are Lipschitz continuous and correctly specified, the sensitivity parameter $\gamma'$ is correctly specified, and the estimating equations for the Q-functions each have a unique solution, then under certain additional regularity conditions, 
% \begin{align*}
%     \widehat{\bm{\theta}}_{1,SA} - \bm{\theta}_{1}^* &\stackrel{d}{\longrightarrow} N(\bm{0}, \Sigma'_{\bm{\theta}_1}) , \\
%     \widehat{d}_{1,SA}^{\mathrm{opt}}(\bm{h}_1) & \stackrel{p}{\longrightarrow} d_1^{\mathrm{opt}}(\bm{h}_1) \text{ for all } \bm{h}_1 \in \mathcal{H}_1,
% \end{align*}
% where $\bm{\theta}_{1}^*$ stands or the true values of $\bm{\theta}_{1}$. 
% $\widehat{\bm{\theta}}_{1,SA}$ and $\widehat{d}_{1,SA}^{\mathrm{opt}}(\cdot)$ denote the estimators of $\bm{\theta}_{1}$ and $d_1^{\mathrm{opt}}(\cdot)$ with the working model \eqref{missingmodel2} estimated by the SA method, respectively. The regularity conditions and the detailed form of $\Sigma'_{\bm{\theta}_1}$ are given in  Web Appendix  D.
% \end{theorem}

We provide  R Markdown tutorials for implementing the proposed weighted Q-learning approaches  using simulated data in the Supplementary Materials. In Web Appendices B-E, we establish the consistency and asymptotic normality of the estimators in Section~\ref{steps} in the 2-stage scenario.
The form of the asymptotic covariance matrix of the estimated parameters at stage 1 is complicated. 
%Therefore, we suggest using the bootstrap method to estimate it in practice.
Therefore, following \cite{Shao_Wang_2016}, we suggest using the bootstrap method to estimate variances and construct confidence intervals of the parameters in practice. Additionally, when 
%$P\{q_{2,1}(\bm{h}_2;\bm{\psi}_2) =0 \} > 0$ and 
non-regularity of stage 1 parameter estimators is of concern, the $m$-out-of-$n$ bootstrap should be used for the inference (see details in Web Appendix G).


\section{Simulation}

We conducted four simulation studies to evaluate the finite-sample performance of the proposed weighted Q-learning method with conditional missingness probabilities of pseudo-outcomes estimated by  the EE or the SA method, which we denote as the WQ-EE method or the WQ-SA method, respectively. 
In Simulation 1, we focused on the scenario where Assumptions 1-5 were satisfied. In Simulation 2, we evaluated the robustness of the WQ-EE method against the violations of Assumption 5 when (1) the chosen instrumental variable was weakly associated with the missingness indicator and (2)  there were interactions between $\bm{u}_t$  and $y_{pse,t}$ in the true conditional missingness probability model. 
In Simulation 3, we demonstrated our simulation-based approach to calibrating the  sensitivity parameter and the performance of the WQ-SA method based on the specified sensitivity parameter values. In Simulation 4, we considered a 3-stage setting where the data generation mechanism  and  the true Q-functions in the first two stages were the same as those in Simulation 1. By comparing the results in Simulations 1 and 4, we investigated the impact of an increased number of total stages on the performance of the proposed methods.
For comparison, we also assessed three other methods for handling partially missing covariate data in Q-learning: a naive method which ignored partially observed covariates (`naive'), complete case analysis (`CC'), and multiple imputation (`MI') \citep{Little_Rubin_2014}. Q-learning when all the covariates are fully observed (`All') was used as the benchmark. More details % of the simulation studies 
can be found in Web Appendix H.

All methods were evaluated in two aspects: Firstly, we compared their performances for estimating the blip function parameters $\bm{\psi}_t$. Secondly, we assessed their abilities to identify the true optimal DTR and compared their values of the estimated optimal DTRs. In Simulation 1, the WQ-EE method and CC analysis provided consistent estimates of blip function parameters at stage 2 (i.e., the final stage) due to the future-independent missingness assumption being satisfied and fully observed pseudo-outcomes. In contrast, naive and MI methods resulted in significant biases at this stage due to the misspecified Q-function models that excluded partially missing covariates in the naive method and the violation of ignorable missingness assumption required by the MI method, respectively.
%\tb{For the naive method, this was because the Q-function models were misspecified with partially missing covariates excluded. For the MI method, this was due to the ignorable missingness assumption for ensuring  consistency was violated.}. 
At stage 1, WQ-EE estimators remained consistent, while the CC method showed substantial biases that did not decrease with increasing sample size. This was because, at stage 1, not only were there nonignorable missing covariates, but the pseudo-outcomes were also nonignorably missing.
%Consequently, the ignorable missingness assumption for the pseudo-outcomes, which is crucial for ensuring the consistency of the CC method, was violated.} 
The WQ-EE method demonstrated the highest proportion of correctly identifying the optimal treatments in both stages. In contrast, the DTRs obtained from CC and MI methods had much lower correct classification rates, and the mean final outcomes under the estimated DTRs were lower than that from the WQ-EE method. In Simulation 2, the WQ-EE method demonstrated robustness to mild model misspecification and consistently outperformed the other three methods, yielding higher mean final outcomes and correct classification rates when Assumption 5 was violated. %\todo{need to give a summary of the results in Simulation 3 and 4.}
In Simulation 3, the proposed simulation-based approach determined plausible ranges for the sensitivity parameter $\gamma_1'$ that covered the true values of $\gamma_1'$ across simulations, and the WQ-SA method with varied values of $\gamma_1'$ was robust and outperformed the naive, CC, and MI methods. In Simulation 4, %we designed a 3-stage setting where the data generating mechanism and the true Q-functions in the first two stages were the same with those in Simulation 1. 
 we compared the values and  overall correct classification rates of the estimated DTRs from Simulations 1 and 4. These performance metrics decreased with the increased number of  total stages,  confirming that the impact of nonignorable missing covariates  propagated over time.%and  the estimation at earlier stages was exacerbated  with an increased number of total stages. 
 

\section{Application to the MIMIC-III data}

We applied the proposed methods to the MIMIC-III data \citep{Johnson_et_al_2016} for investigating the optimal 2-stage fluid strategy to treat sepsis patients in ICUs. 

Sepsis is a life-threatening syndrome caused by the body’s response to infection, resulting in damage to its own tissues and organs. Timely and effective fluid resuscitation is crucial for stabilizing tissue hypoperfusion or septic shock in sepsis patients. The `Surviving Sepsis Campaign' guideline strongly recommends administrating at least 30 mL/kg of intravenous (IV) fluid within the first 3 hours of ICU admission \citep{Rhodes_et_al_2017}.  To address the lack of robust evidence for the optimal fluid resuscitation strategy in the early hours of treatment, \cite{Speth_et_al_2022} utilized MIMIC-III data to assess the optimal DTR during the 0-3 hours and 3-24 hours following admission to the medical ICU. Building upon \cite{Speth_et_al_2022} and in line with the recommendation of the `Surviving Sepsis Campaign' that additional fluid administration after initial resuscitation should be guided by reassessing the patient's hemodynamic status, we incorporated hemodynamic variables in determining the optimal fluid resuscitation strategy during the 3-24 hours period after ICU admission (stage 2). 
%\tb{Note that we used the time since medical ICU admission as the time scale, therefore we will suppress `post-ICU admission' when describing the time points for outcome evaluation. }


We followed similar patient selection criteria as outlined by \cite{Speth_et_al_2022}, focusing on adult septic patients admitted to the medical ICU after initially presenting to the emergency department. The detailed cohort eligibility criteria can be found in Web Appendix I. The treatments in both stages, the baseline covariates, and the final outcome were defined according to \cite{Speth_et_al_2022}. Specifically, the stage 1 and stage 2 treatments were categorized as either fluid restrictive ($<$30 ml/kg) or fluid liberal ($\geq$ 30 ml/kg) strategies in 0-3 hours and 3-24 hours post-ICU admission, respectively. The baseline covariates included gender, age, weight, racial groups, and Elixhauser comorbidity score. The outcome of interest was the Sequential Organ Failure Assessment (SOFA) score evaluated at 24 hours post-admission. The SOFA score is a clinical tool used to assess the severity of organ dysfunction by assigning scores to various organ systems. A higher SOFA score indicates a more severe impairment. In our analysis, we used the negative value of the SOFA score at 24 hours post-admission as the final outcome to be maximized by the optimal fluid resuscitation strategy.

For the intermediate variables considered before stage 2 treatment, we included mechanical ventilation and vasopressors within the first 3-hour period, as well as the patients' SOFA score evaluated at 3 hours post-admission. Additionally, we also considered heart rate, blood pressure, Saturation of Peripheral Oxygen (SpO2), respiratory rate, temperature at 3 hours post-admission and urine output within 0-3 hours after admission. These hemodynamic variables contained missing values, which were likely to be nonignorable due to the informative monitoring of the patients in the MIMIC-III database, as discussed previously.
%because patients who are less sick and not prioritized are more likely to have missing values in their medical records. 
The missing proportions of the hemodynamic variables are provided in Web Table 10. 


There were 973 patients in the selected cohort, among them $53.2\%$ were male and $78.7\%$ were Caucasian. 45.5\% of the patients received the fluid liberal strategy in 0-3 hours post-admission while $53.5\%$ received the fluid liberal strategy in 3-24 hours post-admission. Overall, 67.1\% of patients had fully observed covariates at stage 2, while the rest of the patients had at least one covariate with missing values. 


As discussed in Section 2.3, it was plausible that Assumption 4 holds, that is, the missing data mechanism of the hemodynamic variables at 3 hours post-admission was unrelated to the patient’s SOFA score at 24 hours post-admission, conditioning on the patient's medical history and treatments up to stage 2. Denote the negative of SOFA score at 24 hours post-admission under the optimal stage 2 treatment as $Y_{pse,1}$. It was plausible that weight and gender were independent of the missingness of hemodynamic variables given baseline covariates and $Y_{pse,1}$. Since weight and gender were associated with $Y_{pse,1}$,  they were suitable candidates for nonresponse instrumental variables. Additionally, the Elixhauser comorbidity score can be regarded as a mismeasured proxy of $Y_{pse,1}$, which may also serve as a candidate for the nonresponse instrumental variable,  according to \cite{Miao_Tchetgen_2018}. We investigated the impact of using different sets of candidates for nonresponse instrumental variables in Table~\ref{RDAtable1}. Besides, we applied CC, MI, and the naive method that ignored hemodynamic variables for comparison. The improvement of SOFA at 24 hours was then calculated as the difference between the means of the observed  SOFA scores at 24 hours and the expected  SOFA scores at 24 hours under the estimated optimal DTR. 
%We performed 1000 bootstrap resamplings and applied the 5-fold cross-validation procedure to each bootstrap sample to assess the performance of the estimated optimal DTRs. More details of the estimation procedure are given in Web Appendix H.
We performed cross-validation and $m$-out-of-$n$ bootstrap to assess the performance of the estimated optimal DTRs. More details of the estimation procedure are given in Web Appendix I.
% \tbrown{Several points about table 1 and 3 in the MIMIC data analysis.\\
% \begin{enumerate}
%     \item The choice of $\alpha$ is influenced by the random seed a lot, actually, even when $\alpha = 0$, in most cases the coverage probability would be larger than 0.92 (which is the lower bound for CI when $B_1=200$). So in the end I have arranged the random seed and select $\alpha$ as 0.025. MI and Naive have slightly larger sample sizes because they have lower variance for stage-2 parameter estimators (MI, larger sample size; Naive, less variables) and thus lower $\widehat{p}_{n-regu}$
%     \item I am wondering whether we should change the calculation of the mean and use the original data instead of bootstrap for cross validation (just follow \cite{xu2015regularized}). I have not found paper talked about this issue. But in the real data part of \cite{chakraborty2010inference,chakraborty2013inference,xu2015regularized}, they all use the original data to calculate the point estimates and $m$-out-of-$n$ bootstrap for CI. Although the reviewer did not talked about this issue this time, I think he might ask this next time (Afterall, he did not mention the non-regularity issue last time while just talked about the coverage probability)
% \end{enumerate}
% }  

% \tb{so for 1), the double bootstrap did not pick up large $\alpha$, which means non-regularity did not affect the coverage much in this example, right?   for 2), I think it is fine to use original data for point estimation. I feel like this is some subtle difference between using the traditional statistical methods and using modern prediction models (so with cross-validation). We have demonstrated the improved classification rates using simulations. In the example, we don't really know the optimal DTR, so maybe we can just use the original data without cross-validation? or you meant cross-validation without bootstrap based on original data? I remember in the discussion board, we agreed that we will use original data with repeated sample splitting for point estimation. }

% \tbrown{1, Well, the estimated $\widehat{p}_{n-regu}$ is almost 1 in the MIMIC data analysis part, so a small $\alpha$ can also led to a significant decrease in sample size. If we compare $m/n$ instead of $\alpha$ in real data analysis and simulation, we would find $m/n$ are all around 0.8-0.9.\\
% 2. I mean the cross-validation without bootstrap based on original data. Just as you say,  the original data with repeated sample splitting for train set and test set. I will change to this way now and update the Table 1 and 3 and the description in Web Appendix.} \tb{I see. All look good to me.}

\begin{table}[htbp]
\caption{Improvement of SOFA at 24 hours post-admission for sepsis patients in the selected cohort of  the MIMIC-III database under the estimated optimal fluid strategies with different approaches to handling nonignorable missing covariates.}\label{RDAtable1}
  \centering
  \begin{threeparttable}
    \begin{tabular}{llccc}
    \hline
     & Mean & $m$ & Bootstrap std & Bootstrap 95\% CI \\
    \hline 
    % WQ-EE (ECS)    & 1.294 & X & 0.163 & (0.971, 1.634) \\
    % WQ-EE (weight)   & 1.439 & X & 0.156 &  (1.141, 1.760) \\
    % WQ-EE (ECS, weight)   & 1.334 & X & 0.157 & (1.032, 1.668) \\
    % WQ-EE (gender)   & 1.428 & X&  0.181 & (1.113, 1.775) \\
    % WQ-EE (gender, ECS)   & 1.311 & X & 0.167 &  (1.052, 1.716)\\
    % WQ-EE (gender, weight)    & 1.462 & X & 0.165 & (1.167, 1.799) \\
    % WQ-EE (gender, weight, ECS)    & 1.339 & X & 0.170 & (1.044, 1.692)  \\
    % CC &  1.193 & X&  0.148  &  (0.924, 1.476)     \\
    % MI &  1.093 & X&  0.110  &   (0.884, 1.316)     \\
    % naive & 0.911 & X&0.102   & (0.785, 1.127)  \\
    WQ-EE (ECS)    &  1.297  & 823 & 0.186 & (0.929, 1.667) \\
    WQ-EE (weight)   & 1.432  & 823 & 0.179 &  (1.083, 1.783) \\
    WQ-EE (ECS, weight)   & 1.338 & 823 & 0.181 & (0.976, 1.681) \\
    WQ-EE (gender)   & 1.430 & 823 &  0.210 & (1.012, 1.841) \\
    WQ-EE (gender, ECS)   & 1.338  & 823 & 0.194 &  (0.930, 1.696)\\
    WQ-EE (gender, weight)    & 1.459 & 823 & 0.193 & (1.083, 1.833) \\
    WQ-EE (gender, weight, ECS)    & 1.345 & 823 & 0.197 & (0.949, 1.730)  \\
    CC &  1.190  & 823 &  0.168  &  (0.861, 1.523)     \\
    MI &  1.092  & 827 &  0.123  &  (0.853, 1.340)     \\
    naive &  0.909 & 831 & 0.116   & (0.685, 1.135)  \\
    \hline
    \end{tabular}%
\begin{tablenotes}
    \small
    \item Note: 
    %Mean, the mean of bootstrap point estimates under 5-fold cross-validation; Bootstrap Std, the sample standard error of 1000 cross-validation estimates with bootstrap samples;  Bootstrap 95\% CI, 95\% confidence interval constructed by bootstrap percentiles; ECS, Elixhauser comorbidity score; WQ-EE ($\cdot$), the WQ-EE method based on the corresponding nonresponse instrumental variables listed in parenthesis.
    %Mean, the mean of \tb{$n$-out-of-$n$} bootstrap point estimates under 5-fold cross-validation; 
    Mean, the average of estimated improvement  of SOFA based on cross-validation using 1000 random splits;
    $m$, the  resample size in the $m$-out-of-$n$ bootstrap selected using the double bootstrap method described in Web Appendix G;
    Bootstrap Std, the standard deviation of the $m$-out-of-$n$ bootstrap estimates;
    Bootstrap 95\% CI, 95\% confidence interval constructed by $m$-out-of-$n$ bootstrap percentiles;
    ECS, Elixhauser comorbidity score; WQ-EE ($\cdot$), the WQ-EE method based on the corresponding nonresponse instrumental variables listed in parenthesis.
\end{tablenotes}
    \end{threeparttable}
\end{table}%



Table \ref{RDAtable1} presents the \iffalse bootstrap \fi estimated mean, standard error, and 95\% confidence interval of the improvement of SOFA score at 24 hours post-admission. Notably, all three methods that considered hemodynamic variables outperformed the naive method, which aligns with the recommendation from the Surviving Sepsis Campaign.
Furthermore, the estimated improvement of SOFA at 24 hours post-admission based on the WQ-EE method were notably larger than those based on the CC and MI methods. Importantly, while the WQ-EE method exhibited robustness to the choice of the three candidate nonresponse instrumental variables, the estimated improvement of SOFA at 24 hours post-admission would be lower if the Elixhauser comorbidity score was included in the set of  nonresponse instrumental variables. This could be attributed to the association between the Elixhauser comorbidity score and the missingness of hemodynamic variables. In Simulation 2, we have demonstrated that this association can affect the performance of the WQ-EE method. Therefore, we recommend excluding the Elixhauser comorbidity score from the set of nonresponse instrumental variables. The remaining three combinations of nonresponse instrumental variables showed similar means and standard errors for the estimated improvement of the final outcome. %, and we consider them appropriate nonresponse instrument variables for our investigation.


However, we cannot rule out that gender and weight could be weakly associated with the missingness of hemodynamic variables, even after conditioning on other covariates and the pseudo-outcome. Therefore, we also applied the proposed  WQ-SA method  in case none of the aforementioned variables were valid nonresponse instrument variables. We assumed $\gamma'_1 \geq 0$ because patients with higher SOFA scores (i.e., lower final outcomes) were more severely ill, and they may undergo more intensive monitoring of their hemodynamic status. 
%We considered a range of values for $\gamma'_1$ within the interval $[-2,0]$. The rationale for selecting this range was that when we included gender and weight as nonresponse instrumental variables, the estimated value of $\gamma_1$ in the working model~\eqref{missingnessmodel} results in a bootstrap mean of $-0.630$ with a bootstrap standard error of $0.469$. We believe that the coefficient of $Y_{pse,1}$ would not change considerably if gender and weight were included in the working model~\eqref{missingmodel2} for the conditional missingness probability. Furthermore, when $\gamma'_1=-2$, it deviated by almost three times the standard error from the bootstrap mean, allowing a wide range for the residual association with $Y_{pse,1}$. Table \ref{RDA_sen} shows that when $\gamma'_1=0$, the estimated improvement in SOFA at 24 hours post-admission was nearly identical to that of the CC method, which also assumed a future-independent missingness assumption and ignored the relationship between $Y_{pse,1}$ and $R_{pse,1}$. The estimated improvement in SOFA at 24 hours post-admission increased rapidly as $\gamma'_1$ changed from $0$ to $-1$. When $\gamma'_1$ changes from -1 to -2, the estimated improvement in SOFA at 24 hours post-admission stabilized at around 1.5. 
To calibrate the magnitude of $\gamma'_1$, we applied our simulation-based approach described in Web Appendix A.     Table \ref{RDA_sen_spec} presented the medians of 1000   Wilcoxon rank-sum test p-values for checking the similarity between the `observed' pseudo-outcomes estimates and the \textit{replications} of `observed' pseudo-outcomes simulated under the working missingness model ~\eqref{missingmodel2} with a specified value of $\gamma'_1$. A large median of the p-values indicated that 
the replications of  observed pseudo-outcomes were similar to the observed  pseudo-outcome estimates in the MIMIC-III data, thus 
the corresponding value of $\gamma'_1$ was considered to be plausible. From Table \ref{RDA_sen_spec}, we can  see that the medians of   the p-values were larger than 0.05 when $\gamma_t'$ lied in the range of $[0,1]$ and it was smaller than 0.05 when $\gamma_t' \geq 1.25$. For the sake of prudence, we therefore considered a range of values for $\gamma'_1$ within the interval $[0, 1.25]$.
\begin{table}[htbp]
    \centering
    \caption{Results for the simulation-based approach to determining  a plausible range of the sensitivity parameter for the WQ-SA method  in the  MIMIC-III data application. }
    \begin{tabular}{lccccccccc}
    \hline
    %& \multicolumn{8}{c}{$\gamma_1'$}\\
   $\gamma_1'$  & 0 & 0.25 & 0.5 & 0.75 & 1 & 1.25 & 1.5  & 1.75  \\
    \hline
     Median of  p-values & 0.207 & 0.230 & 0.190 & 0.129 & 0.063 & 0.039 & 0.021 & 0.016  \\
    \hline
    \end{tabular}
    \label{RDA_sen_spec}
\end{table}


Table \ref{RDA_sen} shows that when $\gamma'_1=0$, the estimated improvement in SOFA at 24 hours post-admission was nearly identical to that of the CC method, which also assumed the future-independent missingness assumption and ignored the relationship between $Y_{pse,1}$ and $R_{pse,1}$. The estimated improvement in SOFA at 24 hours post-admission increased rapidly as $\gamma'_1$ changed from $0$ to $0.75$. When $\gamma'_1$ changes from $0.75$ to $1.25$, the estimated improvement in SOFA at 24 hours post-admission stabilized at around 1.47.   
\begin{table}[htbp]
\centering
  \caption{The improvement of SOFA at 24 hours post-admission for sepsis patients in the selected cohort of  the MIMIC-III database under the estimated optimal fluid strategies  with different values of the  sensitivity parameter  $\gamma_1'$ in the WQ-SA method}
  \centering
  \begin{threeparttable}
  \label{RDA_sen}
    \centering
    \begin{tabular}{rcccc}
    \hline
    $\gamma'_1$ & Mean & $m$ &  Bootstrap std & Bootstrap 95\% CI \\
    \hline
    % 0 & 1.187 & \tb{823} & 0.155 & (0.886, 1.485) \\
    % 0.25 & 1.285 & \tb{823} &  0.158 & (0.991, 1.587) \\
    % 0.5 & 1.365 & \tb{823} & 0.160 & (1.052, 1.677) \\
    % 0.75 & 1.432 & \tb{823} & 0.172 & (1.103, 1.782) \\
    % 1 & 1.477 & \tb{823} & 0.174 & (1.147, 1.818) \\
    % 1.25 & 1.498 & \tb{823} & 0.177 & (1.153, 1.885) \\
    0 & 1.184  & 823 & 0.178 & (0.837, 1.534) \\
    0.25 & 1.288  & 823 &  0.181 & (0.930, 1.642) \\
    0.5 & 1.369  & 823 & 0.185 & (1.003, 1.726) \\
    0.75 & 1.436  & 823 & 0.198 & (1.044, 1.821) \\
    1 & 1.475  & 823 & 0.201 & (1.082, 1.873) \\
    1.25 & 1.486  & 823 & 0.204 & (1.085, 1.887) \\
    \hline 
    \end{tabular} 
    \begin{tablenotes}
    \small
    \item Note: 
    %Mean, the mean of bootstrap point estimates under 5-fold cross-validation; Bootstrap Std, the sample standard error of 1000 cross-validation estimates with bootstrap samples;  Bootstrap 95\% CI, 95\% confidence interval constructed by bootstrap percentiles.
    %Mean, the mean of \tb{$n$-out-of-$n$} bootstrap point estimates under 5-fold cross-validation; 
    Mean, the average of estimated improvement  of SOFA based on cross-validation using 1000 random splits;
    $m$, the  resample size in the $m$-out-of-$n$ bootstrap selected with the double bootstrap method described in Web Appendix G;
    Bootstrap Std, the  standard deviation of the $m$-out-of-$n$ bootstrap estimates;
    Bootstrap 95\% CI, 95\% confidence interval constructed by $m$-out-of-$n$ bootstrap percentiles.
    \end{tablenotes}
\end{threeparttable}   
\end{table}%


% \tp{Not so sure how to explain why the estimated SOFA improvement would become larger when $\gamma_1'$ get lesser...\\  
% An update, I actually use the mean of $Y_{pse,1}$ in all sample as the estimated final outcome expectation, not weighting the $Y_{pse,1}$ in the $R_1=1$ sample.\\
% One potential reason for the increase of the estimated improvement might be: the $Y_{pse,1}$ is calculated by |coefficients*covariates| + coefficients*covariates, and when the absolute value of $\gamma_1'$ gets larger, the variance of the coefficients would get larger (in simulation, the parameter estimates from WQ-EE and WQ-SA methods have larger variance than those from CC method), leading to larger absolute value of the coefficients (and we can see that the BS Std gets larger in Table \ref{RDA_sen}). So the |coefficients*covariates| term in calculating $Y_{pse,t}$ tend to be larger when the absolute value of $\gamma_t'$ increase. }


While the challenge of optimizing the management of septic patients is intricate and multifaceted, we employed the proposed methods to investigate whether incorporating partially missing hemodynamic variables could influence the fluid administration after the initial resuscitation to enhance the overall outcome of the septic patients. In line with the best practice recommendations of the Surviving Sepsis Campaign, our analysis results indicated that certain hemodynamic variables should be taken into account when formulating the fluid resuscitation strategy during the 3-24 hours following admission to the medical ICU because accounting for these nonignorable missing covariates with the proposed methods increased the estimated improvement of SOFA at 24 hours post-admission.


\section{Discussion}

In this article, we proposed two weighted Q-learning approaches to estimating optimal DTRs with nonignorable missing covariates and consequent nonignorable missing pseudo-outcomes. To the best of our knowledge, this is the first attempt to estimate optimal DTRs in these challenging nonignorable missingness scenarios. Specifically, we utilized the future-independent missingness assumption and developed a semi-parametric estimating equation-based approach to estimating the conditional missingness probability of the pseudo-outcomes with valid nonresponse instruments. After obtaining the estimated conditional  missingness probabilities, we applied inverse probability weighting to estimate the Q-function parameters. When nonresponse instrumental variables were not available, we developed a practical  approach to calibrating the sensitivity parameter in the WQ-SA method and conducted a sensitivity analysis to assess the impact of varied sensitivity parameter values on the estimated optimal DTRs.  
Our simulation results showed that the WQ-EE and WQ-SA estimators were consistent when their corresponding assumptions hold. They were also robust and outperformed the competing methods under mild model misspecifications. 
In the application to the MIMIC-III data,  we showed that incorporating hemodynamic variables into optimal fluid strategy estimation and handling nonignorable missing covariates with the proposed methods can improve the SOFA score for sepsis patients in ICUs.


In practice, the ignorable missingness assumption is unverifiable from observed data. Therefore, when researchers are not certain about missingness assumptions based on substantive knowledge,  we suggest using the proposed WQ-EE or WQ-SA method to accommodate possible departures from ignorable missingness and  the uncertainty brought by non-ignorable missingness when estimating optimal DTRs.  Notably, the proposed methods are still valid under ignorable missingness (i.e., $\gamma_t=0$ in Assumption 5, $\gamma'_t=0$ in Assumption 6).


There are several directions for future work. First, this article exclusively addresses the issue of missing pseudo-outcomes due to MNAR covariates in optimal DTR estimation with Q-learning. However, in observational studies or EMR databases, patients may not progress through all stages of clinical intervention due to experiencing a terminal event or being censored before the end of the study. In such cases, it is desirable to integrate the terminal event and longitudinal outcome measurements for estimating optimal DTRs while accounting for the censoring and MNAR covariates. Second, the dimension of the covariates is often high in practice, especially in multiple-stage settings. In such scenarios, the kernel regression method is susceptible to the `curse of dimensionality' problem. To address this challenge, one may consider dimension reduction techniques proposed in \cite{Tang_Zhao_Zhu_2014}. 
Third, our methods utilize inverse probability weighting on the  sample with complete covariates at different stages, whose size can be small  when the majority of patients have at least one covariate with missing values. % there are many covariates with missing values across patients\fi many covariates are missing and their missingness indicators are not positively correlated. %when the missingness indicators of different covariates are oppositely distributed. 
In this case, we could explore the multiple imputation method under nonignorable missingness \citep{tompsett2018use} as a more practical alternative.  
Last, the proposed approach can be extended to other  DTRs methods, such as A-learning. 









\vspace{-8pt}
\bibliographystyle{biom} 
\bibliography{DTR_MNAR_bib}



\end{document}
